The graphical identity follows by keeping the same global SCM witness.  A mutual witness is a one-sided witness and has \(Q_a(x,\cdot)\ll P_a(x,\cdot)\) by definition.  Conversely, the one-sided completeness theorem gives \(P_a(x,\cdot)\ll Q_a(x,\cdot)\) on a common conull set; adding \(Q_a(x,\cdot)\ll P_a(x,\cdot)\) makes the same witness satisfy conditional mutual absolute continuity.

For the mixture identity, write
\[
 Q_a=e_aP_a+(1-e_a)R_a.
\]
If \(R_a(x,\cdot)\ll P_a(x,\cdot)\), then \(Q_a(x,\cdot)\ll P_a(x,\cdot)\).  Conversely, \((1-e_a(x))R_a(x,\cdot)\le Q_a(x,\cdot)\), and \(1-e_a(x)\ge\kappa>0\); hence \(Q_a(x,\cdot)\ll P_a(x,\cdot)\) implies \(R_a(x,\cdot)\ll P_a(x,\cdot)\).  These implications hold for each arm on the common conull set, proving the two intersections and the automatic forward-support assertion.